\documentclass[10pt]{article}
\usepackage[T1]{fontenc}
\usepackage[utf8]{inputenc}
\usepackage[french]{babel}
\usepackage{amsmath,amssymb,mathrsfs,mathpazo}
\usepackage{geometry}
\geometry{a4paper,margin=2cm}
\usepackage{graphicx,array,tabularx,tabularray,colortbl,multirow,multicol,booktabs,enumitem,pifont,ulem,fancybox,diagbox,cancel}
\usepackage{pstricks,pstricks-add,pst-plot,pst-tree,pst-node,pst-text,pst-func,pst-math,pst-bspline,pst-3dplot,pst-eucl,pst-solides3d}
\let\pscalculate\undefined
\usepackage{pst-bezier}
\usepackage{tikz,pgfplots,tkz-tab}
\usetikzlibrary{arrows,arrows.meta,patterns,calc,positioning,trees,decorations.markings,shapes.geometric,shapes.misc,shapes.arrows,intersections,decorations.pathmorphing,decorations.pathreplacing,backgrounds,shadows,babel}
\pgfplotsset{compat=1.18}
\usepgfplotslibrary{fillbetween}
\usepackage[np]{numprint}
\usepackage{listings,makecell,pas-tableur,twoopt}
\let\barre\undefined
\usepackage[pstricks]{tablvar}
\definecolor{darkgreen}{RGB}{0,100,0}
\providecommand{\red}{\color{red}}
\providecommand{\blue}{\color{blue}}
\providecommand{\green}{\color{green}}
\providecommand{\black}{\color{black}}
\providecommand{\white}{\color{white}}
\providecommand{\gray}{\color{gray}}
\providecommand{\cyan}{\color{cyan}}
\providecommand{\magenta}{\color{magenta}}
\providecommand{\yellow}{\color{yellow}}
\providecommand{\violet}{\color{violet}}
\providecommand{\orange}{\color{orange}}
\providecommand{\euro}{\texteuro}
\providecommand{\EUR}{\texteuro}
\providecommand{\ieme}{\textsuperscript{e}}
\providecommand{\degres}{\ensuremath{^\circ}}
\providecommand{\degree}{\ensuremath{^\circ}}
\providecommand{\vv}[1]{\overrightarrow{#1}}
\providecommand{\Coord}[1]{#1}
\providecommand{\decosix}{}
\providecommand{\starredbullet}{\ensuremath{\bullet}}
\pagestyle{empty}
\setlength{\parindent}{0pt}
\usepackage{adjustbox}\DeclareUnicodeCharacter{2236}{:}\DeclareUnicodeCharacter{2190}{\ensuremath{\leftarrow}}
\providecommand{\R}{\mathbb{R}}
\providecommand{\N}{\mathbb{N}}
\providecommand{\Z}{\mathbb{Z}}
\providecommand{\Q}{\mathbb{Q}}
\providecommand{\D}{\mathbb{D}}
\providecommand{\C}{\mathbb{C}}
\providecommand{\e}{{\mathrm{e}}}
\providecommand{\vect}[1]{\overrightarrow{#1}}
\providecommand{\vectt}[1]{\overrightarrow{\mathrm{#1}}}
\providecommand{\vvt}[1]{\overrightarrow{\mathrm{#1}}}
\providecommand{\np}[2][]{{#2~#1}}
\providecommand{\vv}[1]{{\overrightarrow{#1}}}
\renewcommand{\d}{{\mathrm{d}}}
\providecommand{\degres}{{^\circ}}
\providecommand{\degre}{{^\circ}}
\providecommand{\blue}{\color{blue}}
\providecommand{\red}{\color{red}}
\providecommand{\enskip}{\quad}
\providecommand{\llbracket}{\lbrack\!\lbrack}
\providecommand{\rrbracket}{\rbrack\!\rbrack}
\providecommand{\euro}{{\text{€}}}
\providecommand{\no}{{\text{n°}}}
\providecommand{\interval}[2]{{[#1;#2]}}
\providecommand{\intervFO}[2]{{[#1;#2[}}
\providecommand{\dotfill}{\ldots}

\providecommand{\ds}{\displaystyle}
\providecommand{\textding}[1]{#1}
\providecommand{\Oij}{\(\mathrm{O};\vec{\imath},\vec{\jmath}\)}
\providecommand{\Oijk}{\(\mathrm{O};\vec{\imath},\vec{\jmath},\vec{k}\)}
\providecommand{\Ouv}{\(\mathrm{O};\vec{u},\vec{v}\)}
\providecommand{\intervFF}[2]{\ensuremath{[#1;#2]}}
\renewcommand{\euro}{\ifmmode\text{\texteuro}\else\texteuro\fi}
\DeclareUnicodeCharacter{00B0}{\ifmmode^{\circ}\else\textdegree\fi}
\DeclareUnicodeCharacter{20AC}{\ifmmode\text{\texteuro}\else\texteuro\fi}
% Énoncés : Éducation nationale. Transcription et corrigés : APMEP, Denis Vergès et contributeurs. Découpage et classement : CoopMaths, Sébastien Lozano, Éric Elter et contributeurs. Adaptation de la présentation : Maths Explorer.
% CC BY-SA : https://coopmaths.fr/www/about/
% https://coopmaths.fr/alea/static/bac/2021/tex/bac_2021_06_metropole_8.tex
\begin{document}
\noindent\textit{APMEP / CoopMaths. CC BY-SA. Présentation adaptée par Maths Explorer.}\par\medskip
\textbf{Commun à tous les candidats}

\medskip

On considère la suite $\left(u_n\right)$ définie par : $u_0 = 1$ et, pour tout entier naturel $n$,

\[ u_{n+1}  = \dfrac{4u_n}{u_n + 4}.\]

\smallskip

\begin{enumerate}
\item ~

\parbox{0.7\linewidth}{
La copie d'écran ci-contre présente les valeurs, calculées à l'aide d'un tableur, des
termes de la suite $\left(u_n\right)$ pour $n$ variant de 0 à 12, ainsi que celles du quotient 
$\dfrac{4}{u_n}$,  (avec, pour les valeurs de $u_n$, affichage de deux
chiffres pour les parties décimales).

À l'aide de ces valeurs, conjecturer l'expression de $\dfrac{4}{u_n}$  en fonction de $n$.

Le but de cet exercice est de démontrer cette conjecture (question \textbf{5.}), et d'en déduire la limite de la suite $\left(u_n\right)$ (question \textbf{6.}).}\hfill
\parbox{0.25\linewidth}{
$\begin{array}{|*{3}{c|}}\hline
n&u_n&\dfrac{4}{u_n}\\ \hline
0 &1,00 &4\\ \hline
1 &0,80 &5\\ \hline
2 &0,67 &6\\ \hline
3 &0,57 &7\\ \hline
4 &0,50 &8\\ \hline
5 &0,44 &9\\ \hline
6 &0,40 &10\\ \hline
7 &0,36 &11\\ \hline
8 &0,33 &12\\ \hline
9 &0,31 &13\\ \hline
10 &0,29& 14 \\ \hline
11 &0,27& 15 \\ \hline
12 &0,25 &16\\ \hline
\end{array}$}

\item Démontrer par récurrence que, pour tout entier naturel $n$, on a : $u_n > 0$.
\item Démontrer que la suite $\left(u_n\right)$ est décroissante.
\item Que peut-on conclure des questions \textbf{2.} et \textbf{3.} concernant la suite $\left(u_n\right)$ ?
\item On considère la suite $\left(v_n\right)$ définie pour tout entier naturel $n$ par : $v_n = \dfrac{4}{u_n}$.

Démontrer que $\left(v_n\right)$ est une suite arithmétique. 

Préciser sa raison et son premier terme. 

En déduire, pour tout entier naturel $n$, l'expression de $v_n$ en fonction de $n$.
\item Déterminer, pour tout entier naturel $n$, l'expression de $u_n$ en fonction de $n$.

En déduire la limite de la suite $\left(u_n\right)$.
\end{enumerate}

\bigskip

\end{document}
